% Appendix: exam/paper_a_solutions.md.

% The answers' headings are labelled ansa:q:N.
\renewcommand{\answerprefix}{ansa}

\chapter{Model Answers to Paper A}
\label{ansa:ch:answers}

Model answers for Appendix~\ref{exa:ch:paper}, with marks shown per part.

\textbf{Code questions are answered as checklists} of what a correct answer must contain, not as
listings. Two correct answers will not look alike, syntax is not marked, and a worked listing here
would be a solution to the lab that asks for the same thing.

Every measured number quoted here was taken on the course's own target, and each comes from the
chapter that publishes it: the FIFO sequence in Question~5 is the Cross-check of
Exercise~\ref{c5:ex:9}, the run in Question~7 is one of the ten in \secref{c7:app:a2}, the interrupt
numbers in Question~8 are those of \secref{c8:app:a1}, and the latencies in Question~10 are those of
\secref{c12:app:a8}.

% ==========================================================================================
\answersection{1}{The image, and the licence on it}{12}

\answerpart{a}{3}
Bootloader, kernel, root filesystem, toolchain.

The \textbf{toolchain} is the one that does not ship. What it determines about the other three is
the \textbf{ABI}: the C library and its version are baked into every binary built with it, so a
glibc binary cannot run on a musl root filesystem, and a binary built against a newer glibc will not
run on an older one. One mark for the four names, one for identifying the toolchain, one for naming
the ABI consequence rather than saying ``it compiles them''.

\answerpart{b}{4}
Two thirds of a mark each, rounded in the candidate's favour:

\begin{booktable}{@{}>{\hsize=0.43\hsize}L>{\hsize=1.57\hsize}L@{}}
  \thead{Component} & \thead{Must publish} \\
  \midrule
  Bootloader & Complete corresponding source, including the board port, its config and build instructions \\
  Kernel & The same, including the in-tree driver, the \code{.config} and the device tree sources \\
  BusyBox & Source at the version shipped \\
  MQTT library (Apache-2.0) & Nothing. Carry the licence text and the NOTICE file \\
  Control daemon & Nothing \\
  Out-of-tree driver & Nothing, if kept proprietary \\
\end{booktable}

\noindent The in-tree driver being \emph{inside} the kernel tree is the point of that row: it is
part of the work whose source must be published, and a candidate who lists it separately as ``ours,
so nothing'' has missed it.

\answerpart{c}{2}
No, it changes nothing. \textbf{The obligation runs to whoever receives the binary and is not
conditional on modification.} One mark for ``no'', one for stating the rule. ``We did not change
it'' is the plausible wrong answer and scores zero; it does make compliance trivial, since the
answer is the upstream tarball at the version shipped, and a candidate who says that earns the
second mark.

\answerpart{d}{3}
The daemon is settled by the \textbf{syscall exception}, which the kernel's \code{COPYING} declares
at its head and \file{LICENSES/exceptions/Linux-syscall-note} states: user programs using kernel
services by normal system calls are not derived works. That is a clarification granted by the
copyright holder, not something that falls out of the GPL, and a candidate who says so earns full
credit for that half.

The driver is not settled because it is compiled against kernel headers and linked into the kernel's
address space, so the derivative-work argument that the syscall exception excludes for applications
applies to it directly. It is contested rather than decided; the kernel enforces a version of the
argument mechanically through \code{EXPORT_SYMBOL_GPL} without the question ever being litigated.

One mark for naming the syscall note, by \code{COPYING} or by \code{Linux-syscall-note}, one for the
address-space distinction, one for identifying the second as contested rather than asserting an
answer.

% ==========================================================================================
\answersection{2}{Asking a running kernel a question}{8}

\answerpart{a}{2}
They did not exist. There is no file and no storage: opening the path selects a kernel function, and
reading it runs that function, which formats a string from live kernel variables at that instant.
Closing the file discards it. One mark for ``nowhere'', one for the mechanism.

\answerpart{b}{3}
Any two of, one and a half marks each:
\begin{itemize}
  \item \textbf{Every file appears to have size zero}, because the kernel does not know the length
    until it generates the answer. A program that trusts \code{stat} to size a buffer reads nothing.
  \item \textbf{Two reads give two different answers}, with no consistency guarantee between them. A
    program computing a rate from one read of a counter has computed nothing.
  \item \textbf{Most of it cannot be mapped or sought from the end.} A program that \code{mmap}s
    it, or seeks to the end to find its length, fails.
  \item \textbf{A read can be expensive}, walking kernel structures under a lock. Polling in a tight
    loop is measurably costly.
\end{itemize}

\answerpart{c}{2}
\file{/proc} grew organically and holds whatever anyone put there, so its files share no format.
\file{/sys} is generated from the driver model, with the rule of one value per file.
\textbf{\file{/sys} is the documented ABI}, recorded in \file{Documentation/ABI/}. One mark each.

\answerpart{d}{1}
Major and minor. Major selects the \textbf{driver}; minor selects \textbf{which device of that
driver's}. Both halves needed for the mark.

% ==========================================================================================
\answersection{3}{Configuring a kernel, and what it costs}{10}

\answerpart{a}{3}
\code{CONFIG_GPIO_SYSFS} is \textbf{absent entirely} from the final \code{.config}. Not \code{n},
and not \code{\# CONFIG_GPIO_SYSFS is not set}: the line the candidate added is discarded.

The prompt is \code{bool "..." if EXPERT}, so with \code{EXPERT} off the symbol \textbf{has no
prompt}, and a symbol with no prompt cannot be set by a user or by a fragment. \code{olddefconfig}
gives it its default, which is absent.

One mark for ``not there at all'', two for the prompt-behind-\code{EXPERT} reasoning. A candidate
who says ``it is set to n'' has the common misconception and earns one mark at most.

\answerpart{b}{3}
\code{depends on B} means the symbol is not offered unless B is on: a \textbf{precondition, which is
checked}. \code{select B} means turning this symbol on \textbf{forces B on: an imperative, which is
not checked}.

\textbf{\code{select} is the one that can produce a violating \code{.config}}, because it does not
evaluate the selected symbol's own \code{depends on}. Selecting B when B depends on C and C is off
leaves \code{B=y} with its own requirements unmet, and the failure surfaces as a compile or link
error in unrelated code. One mark per definition, one for identifying \code{select} with the
mechanism.

\answerpart{c}{2}
Yes, it is set, provided its \code{depends on} is met. A hidden prompt means the symbol is
\textbf{not asked about}, not that it is off; it then takes its \code{default}, here \code{y}. One
mark for ``yes'', one for the distinction between ``nobody was asked'' and ``the value is n''.

\answerpart{d}{2}
The kernel boots normally, prints its usual messages, and then \textbf{panics}, because the driver
it needs to read the root filesystem is stored on the root filesystem:

\begin{console}
Kernel panic - not syncing: VFS: Unable to mount root fs on unknown-block(0,0)
\end{console}

\noindent One mark for the circularity, one for the shape of the message. Wording need not be exact;
``unable to mount root fs'' is enough.

% ==========================================================================================
\answersection{4}{A module, written out}{12}

\answerpart{a}{7}
Marked against this checklist. Order and layout are free.

\begin{booktable}{@{}Lr@{}}
  \thead{Element} & \thead{Marks} \\
  \midrule
  An init function returning \code{int}, and an exit function returning \code{void} & 1 \\
  Both registered, with \code{module_init} and \code{module_exit} & 1 \\
  A message on load and a message on unload, using a \code{pr_*} wrapper & 1 \\
  \code{static char *who = "world";} with \code{module_param(who, charp, ...)} & 1 \\
  A mode of \code{0444}, or any read-only mode. \code{0644} is the trap and costs this mark & 1 \\
  \code{MODULE_LICENSE}, plus author and description & 1 \\
  \code{\#define pr_fmt(fmt) KBUILD_MODNAME ": " fmt} \textbf{before} the includes & 1 \\
\end{booktable}

\noindent Deduct nothing for a missing header, a missing \code{__init}/\code{__exit}, or
\code{MODULE_PARM_DESC}. Deduct the mode mark for a writable mode, since the question states the
requirement explicitly.

\answerpart{b}{2}
\textbf{Nothing of the candidate's is running.} The module is \emph{resident and idle}: code and
data occupying kernel memory, waiting to be called by something else. One mark for ``nothing
running'', one for ``resident''. A candidate who says ``the module keeps running'' has the central
misconception of the lecture and earns zero here.

\answerpart{c}{1}
It places the function in a section that is \textbf{discarded once initialisation is complete},
returning that memory to the system; the boot message \code{Freeing unused kernel memory} is this
happening. Naming the section-and-free behaviour earns the mark; ``it marks it as init code'' does
not.

\answerpart{d}{2}
Any two of: the kernel version; \code{SMP}, so a kernel built without SMP has different locking
primitives; the preemption model; \code{mod_unload}; \code{modversions}; the architecture. Half a
mark each.

The mismatch is an error rather than a warning because \textbf{the kernel has no stable internal
ABI}: a structure that gained a field between two releases would be read at the wrong offsets, and
the result would be silent memory corruption rather than a refusal. One mark for turning corruption
into an error message.

% ==========================================================================================
\answersection{5}{The character device contract}{12}

\answerpart{a}{5}
One mark per line. Capacity 256, starting empty.

\begin{narrowtable}{@{}llr@{}}
  \thead{Step} & \thead{Returns} & \thead{Level after} \\
  \midrule
  write 200 & 200 & 200 \\
  write 100 & \textbf{56} & 256 \\
  read 512 & 256 & 0 \\
  read 512 & \code{-EAGAIN} & 0 \\
  write 300 & 256 & 256 \\
\end{narrowtable}

\lecturefigure[0.84]{exam/images/fifo_sequence.png}
  {A grouped bar chart of the five calls, showing what each asked for against what it returned. The
   second write asks for 100 and returns 56, a short write rather than an error; the second read
   returns no bar at all and is marked -EAGAIN, because an empty FIFO is not end of file.}
  {fig:fifo_sequence}

\answerpart{b}{2}
It returns \textbf{56}, a short write. The FIFO is not full when the call arrives, so \code{-ENOSPC}
does not apply; the contract is that a short count is normal, not an error, and the caller is
expected to read the return value. One mark for 56, one for naming the contract.

\code{-ENOSPC} and \code{100} are the two plausible wrong answers.

\answerpart{c}{3}
Three quarters of a mark each:

\begin{booktable}{@{}LlL@{}}
  \thead{Situation} & \thead{Return} & \thead{Userspace sees} \\
  \midrule
  Non-blocking, nothing buffered & \code{-EAGAIN} & \code{read} returns $-1$, \code{errno} is \code{EAGAIN} \\
  Write with the buffer full & \code{-ENOSPC} & $-1$, \code{errno} is \code{ENOSPC} \\
  \code{copy_to_user} fails & \code{-EFAULT} & $-1$, \code{errno} is \code{EFAULT} \\
  A read of zero bytes & 0 & \code{read} returns 0, \code{errno} is unchanged \\
\end{booktable}

\noindent The last is not an error at all. The \code{read} that Chapter~\ref{c5:ch} specifies
returns 0 for a \code{count} of 0 (\secref{c5:app:b2}): nothing was asked for, so nothing is wrong.
\code{-EINVAL} is the plausible wrong answer and scores nothing for that line.

\answerpart{d}{2}
Any three of these four, two thirds of a mark each, capped at 2:
\begin{itemize}
  \item The page may not be present; userspace memory is demand-paged, and a direct dereference
    faults where no handler expects it.
  \item The address may not belong to the caller, and may point into kernel memory, so a driver that
    copies from it hands kernel memory to a program that asked for it.
  \item The mapping can change under you; another thread can \code{munmap} it between check and
    copy.
  \item Faulting it in may sleep, which is legal here and not in every context the code might be
    reused in.
\end{itemize}

% ==========================================================================================
\answersection{6}{An address, translated and mapped}{12}

\answerpart{a}{2}
\textbf{Two cells}: one address and one size. You know because the \textbf{parent} declares
\code{\#address-cells = <1>} and \code{\#size-cells = <1>}; the node itself does not say. One mark
for ``two'', one for ``the parent declares it''. The second mark is the examinable half.

\answerpart{b}{3}
\code{reg}'s first cell is the child address \code{0x0}. The parent's \code{ranges = <0x00 0x00
0xC000000 0x2000000>} maps child \code{0x0} to parent \code{0x0_0c000000}. So:

\[
  0\text{x}0 + 0\text{x}c000000 = 0\text{x}c000000
\]

\noindent The window is \code{0x1000}, 4096 bytes. Two marks for the address with working shown, one
for the size. A candidate who answers \code{0x0} has taken \code{reg} at face value and earns zero
for this part; follow that error through into part~(c) rather than penalising it twice.

\answerpart{c}{3}
Either:
\begin{itemize}
  \item the address is mapped but is not the device, and the identity register reads
    \textbf{plausible zeros} with no error at all; or
  \item the address is not backed by anything, and the read raises a \textbf{synchronous external
    abort}, killing \code{insmod} and leaving the module half-loaded.
\end{itemize}
One mark each. The third mark is for saying that \textbf{the zeros are worse to debug}: the abort is
immediate, loud and points at the faulting instruction, whereas a driver reading zeros carries on
and misbehaves somewhere else later.

\answerpart{d}{2}
Type \code{0} is an \textbf{SPI}; number \code{0x70} is \textbf{112}; flag \code{4} is
\textbf{level-triggered, active high}. The GIC number is $112 + 32 = 144$. One mark for the three
fields, one for 144.

\answerpart{e}{2}
\code{devm_platform_ioremap_resource(pdev, 0)}.

One mark for the call, one for naming what it replaced: the address the driver used to compute by
hand (the OF core translated \code{reg} through every \code{ranges} when it created the device, so
the driver now only asks for resource 0), \code{request_mem_region}, \code{ioremap}, and the
matching releases in the error path and in \code{remove}.

% ==========================================================================================
\answersection{7}{A race, traced}{10}

\answerpart{a}{1}
$4 \times 20{,}000 = 80{,}000$.

\answerpart{b}{2}
$80{,}000 - 63{,}562 = 16{,}438$ lost, which is $16{,}438 / 80{,}000 = 20.5\%$. One mark each.

\answerpart{c}{3}
\textbf{Three}: a load, an add, and a store.

The interleaving: CPU A loads the counter and gets \emph{n}. CPU B loads the counter, also gets
\emph{n}. Both add one and both store \emph{n+1}. Two increments happened and the counter advanced
by one. One mark for ``three'', two for an interleaving that actually loses an update; a candidate
who says ``they happen at the same time'' without the load-load-store-store detail earns one.

\answerpart{d}{2}
\textbf{No, the bug is absent at no size.} What differs is that the threads do not overlap for long
enough to collide: at 200 iterations the work finishes inside the window it takes the other threads
to get going, so the interleaving never occurs.

One mark for ``not absent'', one for the overlap explanation. This is the question that
distinguishes a candidate who has understood the lecture, and the plausible wrong answer, ``the race
needs more iterations to become likely'', earns one mark for being nearly right about the mechanism
and missing that it is about \emph{duration of overlap} rather than \emph{number of attempts}.

\answerpart{e}{2}
For the counter alone, \textbf{yes}: \code{atomic_inc} makes the load-add-store indivisible and
fixes it completely.

With a running maximum alongside, \textbf{no}. The invariant now spans two variables and must hold
across a read-compare-write sequence, which per-variable atomicity does not provide.

The rule: \textbf{atomics protect one variable; locks protect an invariant.} If the property you
need mentions more than one variable, you need a lock. One mark for the pair of answers, one for the
rule.

% ==========================================================================================
\answersection{8}{An interrupt, acknowledged}{10}

\answerpart{a}{3}
One mark each:
\begin{itemize}
  \item The \textbf{device tree cell}, \code{112}, written by whoever described the board.
  \item The \textbf{GIC hardware number}, \code{144}, which is 112 plus the 32-entry SPI base.
  \item The \textbf{Linux IRQ number}, \code{17}, allocated at run time when the mapping is created
    and unrelated to either. It changes if the machine is booted with a different set of devices.
\end{itemize}

\answerpart{b}{4}
Checklist:

\begin{booktable}{@{}Lr@{}}
  \thead{Element} & \thead{Marks} \\
  \midrule
  Reads its own status register first & 1 \\
  Returns \code{IRQ_NONE} when nothing of its own is pending & 1 \\
  Acknowledges by writing \textbf{back the value it read} & 1 \\
  Drains the FIFO until the level reads zero, and returns \code{IRQ_HANDLED} & 1 \\
\end{booktable}

\noindent Deduct the acknowledgement mark for writing \code{0xFFFFFFFF} or an invented mask, since
that clears bits set after the read and loses those events. The order of the acknowledgement and the
drain is not marked here; part~(c) is where it is examined.

\answerpart{c}{2}
An event arriving between the acknowledge and the drain sets the status bit again \emph{and} adds to
the FIFO the handler is about to empty. The handler services it and returns with the status bit
still set, producing \textbf{one spurious interrupt with nothing to do}.

\textbf{No data is lost.} One mark for the sequence, one for stating that this ordering is the safe
one. A candidate who says data is lost has it backwards; that is the \emph{other} ordering.

\answerpart{d}{1}
The line stays asserted, so the handler is re-entered immediately and forever. \code{insmod} never
returns and the machine is unrecoverable. Either ``interrupt storm'' or ``the machine wedges'' earns
the mark.

% ==========================================================================================
\answersection{9}{A reader put to sleep}{8}

\answerpart{a}{3}
The reader evaluates \code{fifo_empty()} and finds it true. \textbf{The interrupt fires here}: the
handler adds data and calls \code{wake_up}, which finds nothing asleep on the queue and does nothing
at all. The reader then sets its state and calls \code{schedule()}, and sleeps waiting for a wakeup
that has already happened: until the next one, and forever if there is none.

Two marks for the interleaving, one for placing the interrupt precisely between the test and the
state change. A candidate who says ``there is a race'' without locating it earns one.

\answerpart{b}{2}
One mark each:
\begin{itemize}
  \item \textbf{Because of the failure in part~(a):} the macro queues the task and sets its state
    \emph{before} testing the condition, so a wake arriving during that window finds the task on the
    queue. Testing a condition rather than a flag is what allows the test to happen after queueing.
  \item \textbf{Because a wake does not mean the condition is true:} several readers may be woken
    and the first to run may take all the data, and spurious wakes happen. The macro loops and
    re-tests.
\end{itemize}

\answerpart{c}{2}
Half a mark each: the bytes and their count; sleep until data arrives; \code{-EAGAIN};
\code{-ERESTARTSYS}.

\textbf{Zero is never correct} and a candidate who offers it for the second or third case loses that
half mark and should have it pointed out.

\answerpart{d}{1}
The \textbf{delay} family (\code{udelay}, \code{mdelay}), which busy-waits and works in any context;
and the \textbf{sleep} family (\code{msleep}, \code{usleep_range}), which gives up the CPU and needs
process context.

The question: \textbf{can this code sleep?} Both halves needed for the mark.

% ==========================================================================================
\answersection{10}{Frameworks, and the trade}{6}

\answerpart{a}{2}
Any three of, two thirds of a mark each: no way for a program that does not know about it to find
it; no documented interface; no existing tools; no power management; no way to express what the
device \emph{is}; an \code{ioctl} ABI the author now maintains indefinitely, including its 32-on-64
compatibility.

\answerpart{b}{2}
It had to implement a \textbf{channel specification} saying the value is a voltage, indexed, with
\code{IIO_CHAN_INFO_RAW}, and a \textbf{\code{read_raw} callback} returning the value.

It did \textbf{not} have to implement a \file{/dev} node, a \code{read()}, an \code{open()}, an
\code{ioctl}, or any sysfs plumbing. One mark each way round.

\answerpart{c}{2}
\textbf{Neither is better in the abstract, and the question cannot be answered without the
requirement.} A throughput-oriented system should prefer \code{PREEMPT}, whose mean is 19\% lower. A
system with a deadline should prefer \code{PREEMPT_RT}, whose worst case is 2.3 times lower, because
a deadline is a statement about the maximum and an average is not evidence about it.

Full marks for an answer that names the requirement as the deciding factor and reads both columns.
One mark for an answer that picks \code{PREEMPT_RT} and defends it only by the maximum without
noticing the mean got worse. \textbf{Zero for an answer that compares the means alone}, which is the
trap.

A candidate may reasonably argue that the difference is within run-to-run noise and that neither
conclusion is safe from one measurement; that earns full marks and is the better answer.
